\chapter{Cohomology and Obstruction Theory}

If one chapter must bear the weight of the title, it is this one. The phrase
``cohomological typing'' should not merely signal mathematical taste. It should explain why
so many parts of the repository feel related. Cohomology is the device by which one records
that local compatibility problems are not noise but structure. In this monograph, that means
that the difference between success and failure is not exhausted by a boolean result. Success
is the existence of a global section; failure is the presence of a cocycle that refuses to be
resolved by the currently chosen cover.

\section*{From local sections to Čech complexes}
Let a cover of the relevant program object be written as \(\mathcal{U} = \{U_i\}_{i \in I}\),
and let \(\mathcal{F}\) be the presheaf of typed local evidence under discussion. Then the
first terms of the Čech complex are
\[
0 \to \Cech^0(\mathcal{U}, \mathcal{F}) \xrightarrow{\delta^0}
\Cech^1(\mathcal{U}, \mathcal{F}) \xrightarrow{\delta^1}
\Cech^2(\mathcal{U}, \mathcal{F}) \to \cdots.
\]
A $0$-cochain is a family of local sections. Its coboundary measures disagreement on pairwise
overlaps. A vanishing coboundary says that the family is compatible; a nontrivial class in
\(\Hc{1}(\mathcal{U},\mathcal{F})\) says that the local data does not glue globally.

For the software situations emphasized in this repository, one often works with covers whose
salient failures are already visible on pairwise overlaps. This is why first cohomology keeps
reappearing. The point is not that higher cohomology is philosophically impossible. The point
is that the repository's main engineering pathologies---mismatched contracts, inconsistent
local semantics, ungluable modules, competing trust claims---already have a natural home in
\(\Hc{1}\).

\section*{Global sections, obstructions, and repair}
The correct moral is therefore short:
\begin{quote}
\emph{A successful verification or synthesis episode is a vanishing first-obstruction claim
relative to a chosen cover and chosen evidence presheaf.}
\end{quote}

This sentence is powerful because it scales. For a plain typing problem, the cover may be a
family of control-flow and data-flow sites, and the presheaf may record refinement
obligations. For a hybrid coherence problem, the cover is enriched by a layer dimension and
the presheaf now records both obligations and trust-bearing witnesses. For a synthesis
problem, the cover is a family of local module obligations. For anti-hallucination, the cover
contains symbol-resolution, contract, and behavior sites. The outer mathematical sentence is
unchanged.

Repair then becomes a map from obstruction classes to refined local structure. In words,
repair is the act of changing the cover, strengthening a local predicate, refining a module
boundary, or introducing a new witness so that the old cocycle becomes a coboundary. This is
why the repository often treats diagnosis and repair as intimately connected. Once the
obstruction has been localized, one already knows where a refinement should be attempted.

\section*{Five manifestations of first cohomology}
The promise of a single project becomes most visible when the same cohomological sentence is
instantiated in several apparently distant workflows.

\paragraph{Bugs.}
In the most classical case, a bug is a non-gluing local fact. A guard site may allow states
that an arithmetic site forbids. A call site may assume a postcondition stronger than what
the callee provides. A proof site may rely on a branch invariant that the code never really
establishes. These disagreements live on overlaps and therefore naturally become $1$-cocycles.

\paragraph{Hallucinations.}
A generated artifact can be structurally well formed yet fail to inhabit the semantic layer,
or it can fail already at the structural layer because it names nonexistent symbols or
grossly misstates a signature. In either case the failure is local-to-global noncoherence.
The artifact does not inhabit the expected hybrid presheaf across the relevant cover. This is
a cohomological failure, not merely a qualitative disappointment.

\paragraph{Equivalence failure.}
Two implementations may agree at many local sites and still fail to produce a global
isomorphism. The equivalence machinery in \repo{src/deppy/equivalence/} becomes much easier
to understand once local semantic correspondences are read as candidate $0$-cochains whose
incompatibilities define a difference cocycle.

\paragraph{Generation failure.}
A module plan is a cover. Local code generation attempts to populate that cover with sections.
If interfaces, assumptions, proof goals, or exported semantics fail to match across overlaps,
one does not merely have ``bad generation.'' One has an obstruction to descent. The agent is
then in the business of refining covers and retrying gluing.

\paragraph{Trust inconsistency.}
Perhaps the most novel case is trust. One local path may deliver Z3-backed evidence, another
only oracle-backed evidence, and a third may contradict both. These are not just numbers on a
dashboard. They are typed local witnesses whose incompatibility can itself be treated as a
non-gluing phenomenon. The language of trust-valued presheaves is therefore not optional if
one wants the system to reason honestly about mixed evidence.

\begin{figure}[h]
\centering
\begin{tikzpicture}[
  box/.style={draw=chapterblue, rounded corners=2pt, fill=chapterblue!6, minimum width=3.2cm, minimum height=0.9cm, align=center},
  obs/.style={draw=accentgold, rounded corners=2pt, fill=accentgold!12, minimum width=3.3cm, minimum height=0.9cm, align=center},
  >=Latex
]
\node[box] (local) {local sections \ typed evidence on sites};
\node[box, right=1.6cm of local] (overlap) {overlap maps \ pairwise compatibility};
\node[obs, right=1.6cm of overlap] (h1) {$\Hc{1}$ class \ unresolved obstruction};
\node[obs, below=1.0cm of local] (bugs) {bugs};
\node[obs, below=1.0cm of overlap] (hall) {hallucinations / equivalence failures};
\node[obs, below=1.0cm of h1] (gen) {generation / trust inconsistencies};
\draw[->, thick, chapterblue] (local) -- (overlap);
\draw[->, thick, chapterblue] (overlap) -- (h1);
\draw[->, thick, accentgold] (h1) -- (bugs);
\draw[->, thick, accentgold] (h1) -- (hall);
\draw[->, thick, accentgold] (h1) -- (gen);
\end{tikzpicture}
\caption{One cohomological sentence, many manifestations. The repository's main failure modes are organized by the same passage from local sections to overlap incompatibilities to first-obstruction classes.}
\end{figure}

\section*{Obstruction mass and the quality of an elaboration}
The repository's newer scalar-valued ideation layer is also easier to justify in this
language. A candidate elaboration should not be judged only by thematic elegance or by an
LLM's enthusiasm. It should also be penalized when it appears likely to generate unresolved
obstruction mass. In the current code, this appears as a penalty term related to the
observed or estimated \(H^1\)-dimension of the ideation result. That move is mathematically
more principled than it first appears. It says that elaborations carrying unresolved gluing
risk should be scored lower even before full generation begins.

In other words, cohomology does double duty. It diagnoses explicit failures after local
claims have been made, and it also helps rank prospective elaborations before the full cost
of construction is paid. This is one of the clearest signs that the project's geometry is
not ornamental. It actively shapes the agent's search.

\section*{Why first cohomology is not only about failure}
There is a temptation to speak as though cohomology entered the story only when something has
gone wrong. That reading is too narrow. The vanishing of first cohomology is not merely the
absence of drama. It is the positive condition that local evidence has achieved global form.
A gluable family of sections is itself a mathematical achievement. It is a certificate that
multiple local accounts of a program agree. For a system like \DepPy{}, which deliberately
keeps many kinds of evidence visible instead of collapsing them, such certificates are the
closest thing to semantic harmony that one can ask for.

\section*{Repair as a research methodology}
This suggests a wider philosophical point. The repository is not only trying to type and
verify programs. It is trying to organize research methodology around the cycle
\[
  \text{observe locally} \to \text{compute incompatibility} \to \text{refine locally} \to \text{glue globally}.
\]
That cycle is visible in program analysis, in theory-pack extension, in ideation, and in
agent-driven synthesis. Once stated this way, it becomes plausible to claim that
``cohomological typing for absolutely everything'' is a methodology and not merely a slogan.

\begin{proposition}
If a subsystem of software engineering can be represented by a cover of local sites, a notion
of typed local evidence, and a repair loop that attempts to turn overlap failures into
compatibilities, then that subsystem is in scope for the repository's cohomological typing
program.
\end{proposition}

The proof is not a theorem internal to one file; it is the architecture of the repository
itself.
